% Supplementary manuscript section.
\section{A single stationary spike at the cubic moment}
\label{sec:stationary-spike}
Keep half noise and put $Y=1/g$. Suppose $\pi=a$ off a set $E$ of
uniform mass $m>0$, and that
\[
 b\le a\le1,\qquad a\le\pi\le M,\qquad M m^{1/3}\le H. \tag{S1}
\]
The constants $b,H$ are fixed across dimensions. The shape of $E$ and
the variation of $\pi$ inside this envelope are arbitrary. Every
two-valued stationary density with a positive floor and bounded cubic
norm satisfies (S1), taking its higher level for $E$. The constant
density gives ordinary noise by uniqueness of positive scaling.

\paragraph{Theorem.}
Under (S1) the spectral gap is at least $c(b,H)>0$, without coordinate
matching. For a two-valued stationary density there are two stronger
conclusions. The ordered spectral error satisfies $\Delta\le C_b\|\pi-1\|_3$
without matching, and its linear order is sharp. If $p>3$
is fixed, $\|\pi\|_p\le B$, and the stationary marginals match through
order $k$, then the ordered spectral error is at most $Ce^{-ck}$,
with constants depending only on $b,B,p$.

\paragraph{Two-step averaging on the spike.}
The floor $\pi\ge a$ and Jensen give $Y\le Pg/a$. Since
$Pg=P^2(\pi Y)$, writing $T=P^2$ gives the pointwise inequality
$Y\le TY+(M/a)T(Y{\bf1}_E)$. The unweighted term has coefficient one.
Iteration and positivity give, for $j=0,1,2$,
\[
 P^jY\le P^jT^nY+(M/a)\sum_{i=1}^nP^jT^i(Y{\bf1}_E). \tag{S2}
\]
As before $\|Y\|_9\le1/b$. Choose
$n=1+\lceil\log_{16}(1+\log(1/m))\rceil$ and $v=1+8\cdot16^n$.
Forward hypercontractivity gives $\|P^jT^nY\|_v\le1/b$, so its
$L^8$ norm on $E$ is at most $m^{1/8-1/v}/b\le e m^{1/8}/b$.
For $i\ge1$, $P^jT^i$ contracts $L^{23/16}$ into $L^8$:
two ordinary noise steps have correlation $1/4$. An input supported on
$E$ gains a factor $m^d$ relative to its $L^8$ norm, with
$d=16/23-1/8>1/3$. For $m$ sufficiently small in terms of $b,H$,
$(M/a)n m^d\le (H/b)n m^{d-1/3}\le1/2$. In (S2) with $j=0$
absorb this term, then apply the result for $j=1,2$. With
$\mu_E=\mu(\cdot\mid E)$ this gives $\|P^jY\|_{8,\mu_E}\le C_b$.

Now $g=P(\pi Y)$ and $Pg=T(\pi Y)$. The spike contributes at most
$CM m^d$ to $\|Pg\|_{8,\mu_E}$ by the same two-step bound. For
$\|g\|_{3,\mu_E}$, use the $L^{3/2}$-to-$L^3$ contraction of $P$ on
$Y{\bf1}_E$ and its conditional $L^8$ bound; this contribution is
at most $CMm^{1/3}$. The baseline contributions are controlled by
(S2), using $\pi Y\le aY+MY{\bf1}_E$. We obtain
\[
 \|Y\|_{8,\mu_E}+\|g\|_{3,\mu_E}+\|Pg\|_{8,\mu_E}\le C_{b,H}.
 \tag{S3}
\]

\paragraph{All subsets of the spike.}
For nonempty $F\subseteq E$ put $\theta=\nu(F)/(Mm)$ and
$r(F)=J(F,F)/\nu(F)$. Then $E_{\mu_E}[{\bf1}_F\pi/M]=\theta$.
Forward hypercontractivity of $Q$ from $L^{3/2}$ to $L^2$ and
H\"older in $\mu_E$ give
\[
 J(F,F)=\|Q(h{\bf1}_F)\|_2^2
 \le M^2m^{4/3}\|Y\|_{8,\mu_E}^2
                 \|{\bf1}_F\pi/M\|_{24/13,\mu_E}^2
 \le C M^2m^{4/3}\theta^{13/12}.
\]
Thus $r(F)\le CH\theta^{1/12}$. If $r(F)\ge1/2$, then $\theta\ge c>0$.
The density of $\nu(\cdot\mid F)$ relative to $\mu_E$ is at most
$1/c$, so (S3) bounds its mean of $g\,Pg$. Also
$E_{\mu_E}P{\bf1}_E\le m^{1/3}$ by the same forward inequality on the
indicator. Hence $E_{\nu(\cdot\mid F)}P{\bf1}_{F^c}\ge1/2$,
reducing the fixed threshold on $m$ if necessary.
Pointwise Cauchy--Schwarz under $P_x$ and the floor give
$K_x(F^c)\ge b[P_x(F^c)]^2/(g(x)Pg(x))$. A second Cauchy--Schwarz
under $\nu(\cdot\mid F)$ yields
\[
 1-r(F)\ge b\,\frac{(E_{\nu(\cdot\mid F)}P{\bf1}_{F^c})^2}
                         {E_{\nu(\cdot\mid F)}[g\,Pg]}\ge c_{b,H}>0 .
\]
In the other branch the exit is already at least $1/2$. The coarea
estimate in (LR1), whose killed-chain part uses only reversibility and
positivity, gives $\|K_{EE}\|\le1-c'_{b,H}$.

\paragraph{The global gap.}
For any cut $D$, the positive Gram multipliers for $D\cap E$ and
$D\setminus E$ add. On $E^c$, $U=aY^2$ has a bounded $L^{9/2}$
norm. Thus
\[
 \|K_{DD}\|\le\|K_{EE}\|+\|U{\bf1}_{D\setminus E}\|_3
             \le 1-c'_{b,H}+C_b\mu(D)^{1/9}.
\]
All sufficiently small cuts have return bounded away from one, at a
threshold independent of $m,L$. Larger cuts of stationary mass at most
one half obey (G3) at $p=3$, as $E_\mu\pi^3\le1+H^3$. If $m$
exceeds the smallness threshold used above, $M$ is bounded in terms of
$b,H$, and the $L^{9/2}$ argument applies on the whole cube. Cheeger's
inequality proves the global gap in both cases.

\paragraph{Full recovery for two stationary values.}
Now $\pi=M$ on $E$. In the small-$m$ regime (S3) gives
$\|U{\bf1}_E\|_3\le C M m^{1/3}$. Write $\delta=\|\pi-1\|_3$. We prove the quantitative bound
$\Delta\le C_b\delta$. As $\Delta\le1$, it suffices to take $\delta$
small in terms of $b$. With $u=E h$ and $V=1-u^2=\|Qh-u\|_2^2$,
Cauchy--Schwarz gives $1=(E hg)^2\le u E(hg^2)=u E(\pi g)$.
Subtracting $u^2$ gives $V\le u\delta\sqrt V$ since
$\|g-u\|_2\le\sqrt V$. Thus $\sqrt V\le\delta$ and $1-u\le\delta^2$.
Applying forward hypercontractivity to $Qh-u$ yields
$\|g-u\|_3\le\delta$ and $\|Pg-u\|_9\le\delta$.
On $\{g\ge b/2\}$ the inverse square is Lipschitz in $g$. On
$A=\{g<b/2\}$, Jensen gives $Pg\ge b/g>2$, so $Y\le 2(Pg-u)/b$
there and $\mu(A)\le\delta^9$. Consequently
\[
 \|Y^2-1\|_3\le C_b(\delta+\delta^2+\delta^3).
\]
If $M\le2$, multiplication by $\pi$ now gives $\|U-\pi\|_3\le C_b\delta$.
If $M>2$, then $Mm^{1/3}\le2(M-1)m^{1/3}\le2\delta$.
Hence $m\le\delta^3$, and for sufficiently small $\delta$ we may use
(S3) with $H=2$. On $E$ both $\|U{\bf1}_E\|_3$ and
$\|\pi{\bf1}_E\|_3$ are $O_b(\delta)$, and off $E$ we have
$\pi=a\le1$. In both cases $\|U-1\|_3\le C_b\delta$.
The Gram comparison $\Delta\le\|QM_{U-1}Q\|\le\|U-1\|_3$
finishes the upper bound.

The linear order is sharp. In dimension two let $\chi=\chi_1\chi_2$
be the parity character. For $0<t<1$ take
$h=(1+t\chi)/\sqrt{1+t^2/4}$. Then
$\pi=1+\epsilon\chi$, with $\epsilon=(5t/4)/(1+t^2/4)$, and
$U=(1+t\chi)/(1+t\chi/4)$. On the span of $\chi_1,\chi_2$,
$Q$ multiplies by $1/\sqrt2$ and $M_U$ has eigenvector
$\chi_1+\chi_2$ with eigenvalue $(1+t)/(1+t/4)$.
Thus $QM_UQ$ has a nonleading eigenvalue
$(1+t)/(2(1+t/4))<1$. As the ordinary second eigenvalue is $1/2$,
\[
 \liminf_{t\downarrow0}\frac{\Delta}{\|\pi-1\|_3}
 \ \ge\ \lim_{t\downarrow0}\frac{3t}{8(1+t/4)\epsilon}
 \ =\ \frac3{10}.
\]
The example obeys every fixed floor $b<1$ for sufficiently small $t$.

For the exponential estimate let $D=\{|g-1|>2^{-k/8}\}$.
The matching-energy argument gives $\mu(D)\le C_B2^{-3k/4}$.
On $D^c$, $\|(U-\pi){\bf1}_{D^c}\|_3\le C_B2^{-k/8}$ for large $k$.
If $m\le e^{-\alpha k}$, then $M m^{1/3}\le Bm^{1/3-1/p}$
is exponentially small, and the $L^{9/2}$ bound on $E^c$ handles
$D\setminus E$. If $m>e^{-\alpha k}$ then $M\le Be^{\alpha k/p}$
and $\|(U-\pi){\bf1}_D\|_3\le C_b M\mu(D)^{1/9}$. Choose
$0<\alpha<p\log(2)/12$. Both bounds decay exponentially; bounded
$k$ is covered by contraction. Finally $QM_{\pi-1}Q$ is exponentially
small by the matched low/high split for $p>3$.

\paragraph{Scope.}
The
\href{https://airr.science/papers/ARR-2026-4T35VQ6MZB8F99F4/versions/v2/}{stationary-moment manuscript}
treats a two-level potential. Here the two values belong to $hPh$;
the reconstructed $h$ may take many values. The
\href{https://airr.science/papers/ARR-2026-2B2W1WQ7DF9N5VW7/versions/v1/}{published shell}
shows why cubic convergence alone cannot yield this conclusion for an
unrestricted stationary law. The argument for the gap permits multiple
heights on the exceptional set under (S1) and controls each of its
subsets. The unrestricted stationary endpoint remains open.
